\section{The coefficient family and analytic comparison}
\label{sec:coefficient-family}
\label{subsec:coefficient-family}

The two analytic evaluations use the same family of coefficient sums.  We
define that family and state the second moment proved by the theta and
reflection arguments, then explain how the moment enters the zero-free
argument.

Retain the field \(F\), the final set \(S\), the ring \(R=\mathcal O_S\),
and the generator convention fixed in \Cref{sec:arithmetic}.  Unless
explicitly stated otherwise, all ideals in this section are integral
exterior ideals.  A roman letter
representing an ideal is always an allowed \(B\)-generator; it need not be
balanced unless this is stated.  In particular, products of allowed prime
generators are permitted.  The row variable \(m\) below ranges over all of
\(R\).

For \(\re s>1\), write
\(\zeta_F^S(s)=\sum_{\mathfrak a}q_{\mathfrak a}^{-s}\) for the Dedekind
zeta function with its factors at \(S\) omitted, and set
\[
 L_F^S(s,\eta)=L^S(s,\eta):=\prod_{\mathfrak p\notin S}
       (1-\eta(\mathfrak p)q_{\mathfrak p}^{-s})^{-1}.
\]
The latter notation also denotes the usual meromorphic continuation of
this incomplete Hecke \(L\)-function; a principal pole is permitted.  All
primes dividing the conductor of \(\eta\) already lie in \(S\), so there
are no further omitted exterior factors for this target character.
For \(t>0\) put
\begin{equation}\label{eq:reflection-gaussian}
 P_G(t)=\frac{1}{2\sqrt\pi}\exp\!\left(-\frac{(\log t)^2}{4}\right),
 \qquad
 \int_0^\infty P_G(t)t^\xi\,\frac{\dd t}{t}=e^{\xi^2}
 \quad(\xi\in\C).
\end{equation}
The transform identity follows by putting \(v=\log t\) and completing the
square, first for real \(\xi\) and then by entire continuation.  Write
\(\vartheta=\overline\Lambda\,\overline G\,\eta\).  For \(m\in R\)
and \(Z>0\) define
\begin{equation}\label{eq:reflection-B-definition}
 B_m(Z)=
 \sum_{\substack{\mathfrak c\ \mathrm{squarefree}\\ \mathfrak n}}
 \frac{\gamma_2(\mathfrak c)}{q_{\mathfrak c}^{1/2}q_{\mathfrak n}}
 \vartheta(\mathfrak A)\chi_{\mathfrak A}(m)
 P_G\!\left(\frac{q_{\mathfrak A}}{Z}\right),
 \qquad
 \mathfrak A=\mathfrak c\mathfrak n^3,\quad A=c n^3.
\end{equation}
Here \(c,n,A\) are allowed generators, and \(\mathfrak n\) is unrestricted,
so it may share primes with \(\mathfrak c\).
Every residue-character power has the zero extension on nonunits fixed in
\Cref{sec:arithmetic}, including powers with exponent zero modulo six.
The sum is independent of the permitted generator choices.  It converges
absolutely: \(|\gamma_2(\mathfrak c)|=1\) for squarefree \(\mathfrak c\),
and for any \(\sigma>1/2\) the bound \(P_G(t)\ll_\sigma t^{-\sigma}\)
reduces the absolute sum, after a factor \(O_\sigma(Z^\sigma)\), to
\begin{equation}\label{eq:reflection-original-dirichlet-majorant}
 \sum_{\mathfrak c,\mathfrak n}
 q_{\mathfrak c}^{-1/2-\sigma}q_{\mathfrak n}^{-1-3\sigma}
 \leq \zeta_F^S(\sigma+1/2)\zeta_F^S(3\sigma+1)<\infty.
\end{equation}
At \(m=0\), the mask kills every term with \(\mathfrak A\ne R\), so
\(B_0(Z)=P_G(1/Z)\).  The factor
\(\gamma_2(\mathfrak c)/(q_{\mathfrak c}^{1/2}q_{\mathfrak n})\) will be
identified in \Cref{sec:theta} as the exterior coefficient supplied by the
normalized local theta formula.

The comparison uses rows at scale \(XY\), where
\begin{equation}\label{eq:reflection-scales}
\begin{gathered}
 a=\frac9{10},\qquad b=\frac{133}{1000},\qquad
 M_0=a+b=\frac{1033}{1000},\\
 h_0=1-a=\frac1{10},\qquad X=Z^a,\quad Y=Z^b.
\end{gathered}
\end{equation}
Recall the norm section \(\ell_T\) fixed in \Cref{sec:arithmetic}, for which
\(|\ell_T|_S=T\).  If \(\mathscr K\subset F_S^\times\) is compact, set
\[
 \mathscr B_Z(\mathscr K)=R\cap\ell_{Z^{M_0}}\mathscr K.
\]
Thus every local absolute value of \(\ell_{Z^{M_0}}^{-1}m\) is bounded
above and away from zero on this set.  These are the compact balanced
annular row sets used below.  For \(m\ne0\) decompose its exterior ideal
uniquely as
\[
 (m)_S=\mathfrak h(m)\mathfrak d(m),
 \qquad
 \mathfrak h(m)=\prod_{v_{\mathfrak p}(m)\geq2}
                    \mathfrak p^{v_{\mathfrak p}(m)},
 \qquad
 \mathfrak d(m)=\prod_{v_{\mathfrak p}(m)=1}\mathfrak p.
\]
The ideal \(\mathfrak d(m)\) is squarefree and coprime to \(\mathfrak h(m)\);
every prime exponent of \(\mathfrak h(m)\) is at least two.  We call such an
ideal powerful.

The theta and reflection arguments in \Cref{sec:theta,sec:reflection}
prove the following estimate.

\begin{proposition}[Reflected second moment]\label{prop:reflection}
Let \(F\) be a fixed cyclotomic field containing \(\mu_{12}\), let \(S\) and the
arithmetic data be as above, and fix a compact subset
\(\mathscr K\subset F_S^\times\).  For every \(\epsilon>0\), every \(Z\geq2\),
and all dyadic lengths \(P,D\geq1\),
\begin{equation}\label{eq:reflection-dyadic-mass}
 \sum_{\substack{m\in\mathscr B_Z(\mathscr K)\\
        P\leq q_{\mathfrak h(m)}<2P\\
        D\leq q_{\mathfrak d(m)}<2D}}
 |B_m(Z)|^2
 \ll_{F,S,\eta,\mathscr K,\epsilon}
 Z^\epsilon P^{1/2}\left(D+\frac{D^2P}{Z}\right).
\end{equation}
In particular,
\begin{equation}\label{eq:reflection-total-mass}
 \sum_{m\in\mathscr B_Z(\mathscr K)}|B_m(Z)|^2
 \ll_{F,S,\eta,\mathscr K,\epsilon} Z^{2M_0-1+\epsilon}.
\end{equation}
The estimate is uniform if \(\mathscr K\) is multiplied by an element of any
fixed compact subset of the norm-one group of \(F_S^\times\).
\end{proposition}

Here is the analytic comparison proved in \Cref{sec:poisson,sec:zero-free}.
The first of those sections constructs a scalar
\(I_Z=I(Z^a,Z^b,Z)\) by pairing \(B_m(Z)\), on a fixed compact annular
row set of scale \(XY\), with a normalized sum of
\(\chi_{\mathfrak s}(d)\e(-md/s)\), where
\(q_{\mathfrak s}\asymp Y\) and \(d\in(R/\mathfrak s)^\times\).
It specifies the weights and the average over \(F_S^1/U_S^6\).
The whole additive factor has squared row mass \(O(Y^{-1})\), uniformly in
that average.  Hence \Cref{eq:reflection-total-mass} and Cauchy--Schwarz give,
for every \(\epsilon>0\),
\[
 I_Z\ll Z^{(2M_0-1-b)/2+\epsilon}
       =Z^{933/2000+\epsilon}.
\]
The constants here and below may depend on the fixed \(F,S,\eta\) and the
fixed weights, as well as on any displayed loss.

Poisson summation gives a second evaluation.  Its nonzero frequencies are
grouped by \(u\in\mathcal U\) through \Cref{eq:frequency-decomposition}.
Only the principal class \(u=1\), the class of sixth powers, contributes
both numerator poles in the auxiliary Mellin variables: \(w=1\) and
\(z=1/6\).  The Euler calculation constructs a specific principal correction
\(\mathcal H_\eta(s)\) by specializing the correction for \(u=1\) at
\((w,z)=(1,1/6)\).  It proves that \(\mathcal H_\eta\) is holomorphic and
nowhere zero on an open neighborhood of the closed half-plane
\(\re s\ge .998\), and is bounded on that closed half-plane uniformly in
height.  The double residue equals \(c_0f_\eta(Z)\), where \(c_0\ne0\) and
\[
 f_\eta(Z)=\frac1{2\pi i}\int_{(2)}
       Z^{s-8/15}e^{(s-5/6)^2}
       \frac{\mathcal H_\eta(s)}{L^S(s,\eta)}\,ds
       \qquad(Z>0).
\]
The exponent left by the two residues is
\(a/2+s-1+h_0/6=s-8/15\), whose value at \(s=1\) is \(7/15\).

The zero-detection argument proves, for sufficiently large \(Z\),
\[
 I_Z=c_0f_\eta(Z)+O(Z^{7/15-.0004}).
\]
Since \(933/2000=7/15-1/6000\), this identity and the preceding bound,
with the stated losses, give \(f_\eta(Z)\ll Z^{7/15-1/20000}\).
Moving the original \(s\)-line in the inverse Mellin integral to the right
gives \(f_\eta(Z)=O_A(Z^A)\) for every \(A>0\) and \(0<Z\le1\).
Consequently the shifted Mellin transform
\(\int_0^\infty f_\eta(Z)Z^{-s+8/15}\,dZ/Z\) is holomorphic for
\(\re s>1-1/20000\).  Fourier inversion on \(\re s=2\), followed by
meromorphic continuation, identifies it throughout that half-plane with
\[
 e^{(s-5/6)^2}\frac{\mathcal H_\eta(s)}{L^S(s,\eta)}.
\]
Its numerator is holomorphic and nowhere zero there, so \(L^S(s,\eta)\)
has no zero there; a principal pole only gives a zero of the reciprocal.
The numerical saving \(1/20000\) is absolute; the constants and thresholds
may depend on \(F,S,\eta\).  The later
sections supply the residue calculation and the contour and Fourier
justifications.
